\chapter{Seeing the Good in Others}

A person's rebirth often begins with a small and embarrassing discovery: he is not as fair-minded as he thought.

He may believe that he respects ability, loves truth, and learns from anyone. Then one day he notices that he cannot bear to see certain people succeed. He does not merely dislike their faults. He refuses to admit their strengths. Their good qualities become invisible to him, not because those qualities are absent, but because his attention refuses to land on them.

This is a serious loss. It is not mainly a moral flaw, though it may become one. It is first an economic error in the deepest sense of this book: a waste of attention.

The road to wealth freedom requires us to collect useful signals wherever they appear. If a person can learn only from those he already likes, respects, envies correctly, or considers ideologically acceptable, then his learning surface is too small. He has voluntarily narrowed the world.

The harder practice is this:

\begin{quote}
Do not let a person's faults blind you to his strengths. Do not let a person's strengths blind you to his faults.
\end{quote}

Both halves matter. Naive admiration is dangerous. So is reflexive contempt. The free person is not the person who likes everyone. The free person is the person whose attention remains available for reality.

\section*{The Shock of Having No Judgment}

One useful way to understand this problem is through intellectual humiliation.

Imagine listening to one person in the morning. He explains something with confidence, structure, and force. You think, ``That makes sense. He is impressive.''

Then, in the afternoon, another person explains something else with equal confidence, structure, and force. Again you think, ``That makes sense. He is impressive.''

That night, while lying in bed, you suddenly realize that the two arguments contradict each other at the root. They cannot both be true in the same way. Yet at the time you accepted both.

The immediate conclusion is painful: perhaps you did not understand either argument. Perhaps you were responding to fluency, tone, and authority rather than truth. Perhaps what you called judgment was only impression.

This kind of shock is valuable. It can push a person toward critical thinking, toward the habit of asking:

\begin{itemize}
\item What is the claim?
\item What assumptions support it?
\item What evidence would weaken it?
\item What alternative explanation exists?
\item Does this conflict with something else I already believe?
\end{itemize}

A person without these questions is easily moved by whoever speaks last, speaks loudest, or speaks most beautifully. He may feel intelligent because he is often impressed, but being impressed is not the same as thinking.

Good books can help. A serious book on logic, judgment, or critical thinking is not merely information; it is a mirror. It allows you to see how often feeling has been pretending to be reasoning. When possible, reading important books in their original language is also worthwhile, because style, precision, and hidden assumptions are often damaged in translation. But the larger point is simpler: when a problem is real, the answer has probably been explored somewhere. Search for the book. Search for the discipline. Search for the tool.

\section*{Four Positions}

One useful model divides human perspective into four positions:

\begin{center}
\begin{tabular}{p{0.38\linewidth}p{0.45\linewidth}}
\hline
Position & Inner sentence \\
\hline
Collapse & I am not OK, and you are not OK. \\
Dependence & I am not OK, and you are OK. \\
Defensiveness & I am OK, and you are not OK. \\
Maturity & I am OK, and you are OK. \\
\hline
\end{tabular}
\end{center}

Childhood often begins near the first two positions. The world is run by adults. Teachers, parents, and older people define the rules. The child is often wrong, small, corrected, delayed, and dependent. Many people carry that early pressure into adulthood.

Then the reaction begins. A person who has long felt ``I am not OK, you are OK'' may try to escape by flipping into ``I am OK, you are not OK.'' This feels like progress, but it is often only revenge in intellectual clothing.

So many adults spend their lives moving between the second and third positions:

\begin{quote}
Either you are better than me, or I am better than you.
\end{quote}

This is exhausting. It makes every interaction a ranking event. Someone else's success becomes your demotion. Someone else's mistake becomes your comfort. Someone else's praise becomes evidence against you. Someone else's weakness becomes evidence for you.

The fourth position is different: ``I am OK, and you are OK.''

This does not mean everyone is equally competent, equally honest, or equally suitable for trust. It means my existence is not threatened by your strength, and your existence is not erased by my strength. It means I can learn from you without kneeling before you. It means I can notice your weakness without needing to destroy you.

This position is rare because it requires attention to stop living inside comparison.

\section*{One Word Changes the Problem}

There is a useful distinction between two sentences:

\begin{quote}
I cannot bear to see others doing well.

I cannot bear to see the good in others.
\end{quote}

The first is about another person's whole condition: success, status, money, reputation, beauty, opportunity, recognition.

The second is more precise. It concerns the good qualities inside another person: discipline, taste, courage, patience, skill, clarity, kindness, timing, execution, or restraint.

A person may have obvious faults and still possess one excellent quality. If your dislike for him prevents you from seeing that quality, then the loss is yours. His strength continues to operate whether or not you admit it. Your blindness does not weaken him. It weakens you.

This is why metacognition matters. We need the ability to observe where attention has gone. When we dislike someone, attention usually rushes toward the evidence that justifies dislike. We notice his errors, exaggerations, hypocrisy, vanity, carelessness, and bad taste. We rehearse them. We collect them. We use them to protect ourselves from the discomfort of admitting that he may also be good at something.

The metacognitive question interrupts the pattern:

\begin{quote}
What is this person doing well that I am refusing to see?
\end{quote}

This question does not require affection. It requires accuracy.

\section*{The Double Standard}

Without metacognition, most people use a double standard.

When looking at others, they attend to what is wrong.

When looking at themselves, they attend to what is right, justified, improving, or misunderstood.

This double standard feels natural because it protects the self. It also blocks growth. If I only study other people's defects, I gain very little. Their defects may warn me, but they rarely enlarge me. If I study their strengths, I can borrow methods, standards, habits, and possibilities.

The useful rule is balanced:

\begin{center}
\begin{tabular}{p{0.35\linewidth}p{0.45\linewidth}}
\hline
What you see & How to use it \\
\hline
Another person's strength & Treat it as material for learning. \\
Another person's weakness & Treat it as a reminder and warning. \\
Your own strength & Build on it without worshiping it. \\
Your own weakness & Repair it without dramatizing it. \\
\hline
\end{tabular}
\end{center}

This is not niceness. It is extraction of value from reality.

If a person you dislike is punctual, learn punctuality. If he writes clearly, study the clarity. If he sells well, understand the sale. If he manages energy well, observe the rhythm. If he takes risk well, examine the preparation behind the risk. If he recovers quickly after failure, notice the mechanism.

You are not required to approve of the whole person in order to learn from one part of him.

\section*{Why It Takes So Long}

The idea is simple. The practice may take years.

Many people understand ``I am OK, and you are OK'' long before they can live there. The mind agrees. The body resists. The old reflex returns under pressure.

Why is it so hard?

One reason is that many people are still trying to prove themselves. Before a person has received any stable recognition, comparison feels unavoidable. He is still asking the world, ``Am I enough? Do I count? Can anyone see me?'' In that condition, another person's good qualities may feel like evidence against him.

This is why financial pressure is so powerful. Money is not the only form of recognition, but lack of money often intensifies the need to compare. When basic survival, dignity, and future security feel uncertain, the mind becomes defensive. It searches for proof that it is not falling behind. It may become addicted to small victories over others.

But comparison has no end.

The beautiful person may lack education. The educated person may lack health. The healthy person may lack background. The well-connected person may lack competence. The competent person may lack courage. The courageous person may lack patience.

A mind devoted to comparison can always find a ranking that hurts or flatters itself. Neither result produces growth.

This is the crucial turn:

\begin{quote}
Put every available unit of attention on your own growth.
\end{quote}

This is one of the foundations of wealth freedom and, in fact, of any serious freedom. The point is not to become self-absorbed. The point is to stop spending attention on theatrical proof.

If your strength is real, proof will eventually become unnecessary. If your strength is not real, performance will not make it real.

\section*{Recognition and Peace}

There is a hard truth here. Some degree of real progress makes generosity easier.

When a person has built something, learned something, earned something, survived something, and gained some confidence through reality rather than fantasy, he has less need to win every invisible contest. He can admire more freely because admiration no longer feels like surrender.

This is another reason to pursue wealth freedom seriously. It is not merely about consumption. It is about reducing the psychological pressure that turns every other person into a rival. When economic anxiety loosens, attention becomes less frantic. When attention becomes less frantic, learning becomes easier. When learning becomes easier, growth accelerates.

Still, waiting for future success is not enough. A person who says, ``I will be generous after I am strong,'' may never begin. The practice must start now, especially while it is difficult.

Use the following small protocol whenever envy, contempt, or comparison appears:

\begin{enumerate}
\item Notice the position: am I feeling ``I am not OK, you are OK'' or ``I am OK, you are not OK''?
\item Name the attention leak: am I studying his faults to comfort myself?
\item Extract one useful strength: what is he doing well?
\item Convert it into practice: what can I imitate, adapt, or study?
\item Return to growth: what is my next concrete action?
\end{enumerate}

This is how metacognition becomes practical. It is not a concept to admire. It is a switch to flip when attention begins to decay.

\section*{Compare Futures}

There will still be moments when comparison seems unavoidable. You meet someone younger, richer, clearer, stronger, calmer, more fluent, more disciplined, more loved, or more prepared. The mind says, ``I am not OK.''

At that moment, compare futures, not current positions.

If you compare only the present, you may become ashamed or arrogant. Both are traps. If you compare futures, the question changes:

\begin{quote}
Given what I now see, what can I become if I start growing from here?
\end{quote}

A person living in the future does not deny the present. He simply refuses to treat it as final. Today's gap becomes information. Another person's excellence becomes evidence that a higher standard exists. The correct response is not resentment. The correct response is design.

This is why seeing the good in others can feel like a rebirth. The same world becomes richer. The people around you have not changed much. Your access to them has changed. Where you once saw threats, you now see models. Where you once saw reasons to complain, you now see warnings. Where you once saw rivals, you now see data.

And occasionally, where you once saw an enemy, you may discover a teacher.

\section*{A Private Exercise}

This chapter asks for an uncomfortable exercise. It does not need to be public. It should be written.

Choose a notebook, or a dedicated section of a notebook, for this one practice: recording the moments when you cannot bear to see the good in others.

Write down the name or situation. Then answer:

\begin{itemize}
\item What exactly bothers me about this person?
\item What good quality am I refusing to admit?
\item Why does that quality threaten me?
\item What can I learn from it?
\item What weakness in me is being exposed?
\item What action would turn this discomfort into growth?
\end{itemize}

Also ask a harder question:

\begin{quote}
What important lesson have I learned from someone I dislike?
\end{quote}

If you can answer this honestly, your attention is becoming freer.

The person you dislike may still be untrustworthy. He may still be wrong in many areas. You may still need distance, boundaries, or refusal. None of that cancels the practice. Accuracy is not approval. Learning is not submission. Seeing clearly is not surrender.

\section*{The Final Shift}

People who cannot see the good in others are often terrified that others will see the bad in them.

This is the hidden symmetry. The person who constantly judges is usually afraid of being judged. The person who erases others' strengths often fears that his own weaknesses will erase him. The person trapped in comparison cannot rest because he imagines everyone else is comparing in the same cruel way.

The exit is not to pretend that weakness does not exist. The exit is to accept a more complete reality:

\begin{quote}
Every person has strengths and weaknesses. So do I.
\end{quote}

From there, life becomes simpler.

Other people's strengths are available for study. Other people's weaknesses are available as warnings. My strengths are available for building. My weaknesses are available for repair.

No drama is required.

The instruction is severe but liberating:

\begin{quote}
Proving yourself is not important. Growth is important. If growth becomes real, proof will take care of itself.
\end{quote}

To see the good in others is not merely to become kinder. It is to recover a large part of the world as usable material. It is to stop paying attention to rank and begin paying attention to reality. It is to build the inner position from which wealth freedom becomes possible: I can grow, and you can grow; I can be good, and you can be good; your excellence does not reduce my future.

The person who reaches that position has not escaped competition by becoming weak. He has escaped useless competition by becoming stronger.
